\section{The Aesthetic Bar: Beating Mermaid Visually}
\label{sec:aesthetics}

Aesthetics is not a finishing coat applied after architecture---it is a \textbf{first-class
pillar} of the product. The explicit goal: every diagram produced by this compiler must look
demonstrably better than its Mermaid equivalent. ``Pro,'' ``neat,'' ``coherent'' are
requirements, not aspirations.

\subsection{Why Aesthetics Is Architecture}

Mermaid is ubiquitous because of its ecosystem integration, not its visual quality.
Users accept its output because it's free and frictionless---not because it looks good.
The aesthetic gap is the primary differentiation lever:

\begin{description}
  \item[Mermaid] Documentation-grade. Fixed styles. Poor typography. No design system.
        Functional but never beautiful.
  \item[This compiler] Presentation-grade. Themeable. Proper typography. Coherent design
        language. Output you'd put in a board deck or technical blog.
\end{description}

This is not subjective polish---it is the \textbf{primary reason users choose this tool
over Mermaid}. If the output doesn't look materially better, there is no product.

\subsection{The Design System}

Every diagram shares a coherent visual language:

\begin{description}
  \item[Typography] System-level font stack with proper weight hierarchy
        (title/heading/body/label/caption). Consistent sizing ratios across all diagram types.
        No all-caps labels. No tiny illegible text.

  \item[Spacing and rhythm] Consistent padding, margins, and gaps derived from a base unit.
        Generous whitespace. Visual breathing room. The ``no clutter'' principle
        (Section~\ref{principle:minimal-clutter}) is enforced structurally.

  \item[Color palette] Curated palettes with proper contrast ratios (WCAG AA minimum).
        Semantic color mapping: status colors, category colors, emphasis colors.
        Dark and light modes as first-class variants, not afterthoughts.

  \item[Shape language] Rounded corners, consistent border weights, soft shadows where
        appropriate. No harsh black-on-white wireframe aesthetic.

  \item[Edge/connector style] Smooth curves, proper arrowheads, consistent stroke widths.
        Routing that avoids overlaps where the layout algorithm permits.
\end{description}

\subsection{Default Theme Craft}

The default theme must be beautiful \emph{out of the box}. Users who never configure a theme
should still produce output that surpasses Mermaid. Design commitments:

\begin{itemize}
  \item Modern, clean aesthetic inspired by technical publications (not wireframes)
  \item Light theme: warm grays, blue accent, Inter/system-ui stack
  \item Dark theme: deep navy/charcoal, blue-teal accent, proper contrast
  \item Professional appearance suitable for slides, blog posts, documentation
  \item No garish colors, no 1990s UML aesthetic, no pixelated edges
\end{itemize}

\subsection{Brand Themes and Coherence}

The theme system enables organization-specific brand themes:

\begin{itemize}
  \item Themes can define complete visual identities (colors, fonts, shapes, spacing)
  \item Single inheritance: branded themes extend base themes, overriding selectively
  \item Cross-diagram coherence: a theme applied to a poster produces visually unified
        output across all constituent diagram types
  \item Built-in exemplars: \texttt{professional-light}, \texttt{professional-dark},
        \texttt{bytebytego} (inspired by the popular technical illustration style)
\end{itemize}

\subsection{Grammar = Semantics; Theme = Style}

The architectural discipline that enables aesthetic excellence:

\begin{quote}
\emph{The Domain IR carries what to communicate. The theme carries how it looks.
Never mix them. A theme change cannot alter diagram meaning. A content change
cannot alter diagram style.}
\end{quote}

This separation means:
\begin{itemize}
  \item The same diagram looks professional in \texttt{professional-light} and playful
        in a hypothetical \texttt{sketch} theme---without any IR changes.
  \item Agents author content without making aesthetic decisions.
  \item Designers craft themes without understanding grammar semantics.
  \item New grammars inherit the full theme system and look coherent from day one.
\end{itemize}

\subsection{Aesthetic Verification}

Beauty is hard to test, but we enforce baselines:

\begin{itemize}
  \item \textbf{Golden-image regression}: visual output is diffed against known-good
        baselines. Any change to rendering requires explicit approval.
  \item \textbf{Contrast ratio checks}: automated WCAG AA verification for all
        text-on-background combinations in built-in themes.
  \item \textbf{Gallery examples}: every grammar ships with gallery outputs demonstrating
        the aesthetic bar. These serve as both documentation and regression targets.
  \item \textbf{Side-by-side comparisons}: marketing/documentation includes Mermaid vs.\
        ours comparisons for every supported diagram type.
\end{itemize}

\subsection{Implementation Status}

\textbf{Built:} The theme system is fully implemented. Five timeline themes (including dark
and ByteByteGo variants), two flow themes, two sequence themes, two tree themes, and two
composition themes are shipped. All enforce the grammar=semantics/theme=style discipline.

\textbf{In progress:} Expanding the aesthetic standard to the new diagram types as they
ship. The UML/software line (Tier-1) will establish the ``modern UML'' visual standard
that differentiates from PlantUML's dated look.
